\documentclass[11pt,letterpaper]{article}
\usepackage[margin=1in,headheight=14pt]{geometry}
\usepackage[T1]{fontenc}
\usepackage{amsmath,amsthm,mathtools}
\usepackage{newtxtext,newtxmath}
\usepackage{microtype,booktabs,array,longtable}
\usepackage[table]{xcolor}
\usepackage{enumitem,fancyhdr}
\usepackage{listings}
\usepackage{hyperref}
\definecolor{ink}{HTML}{17324D}
\definecolor{accent}{HTML}{156B73}
\hypersetup{colorlinks=true,linkcolor=ink,citecolor=accent,urlcolor=accent,
 pdftitle={Finite Fredholm Truncation for Long Increasing Subsequences},
 pdfauthor={ChatGPT},pdfsubject={Fixed-sector holonomicity, OEIS A047874 and A269021}}
\pagestyle{fancy}\fancyhf{}
\fancyhead[L]{\small\textsc{Long increasing subsequences}}
\fancyhead[R]{\small Research article}
\fancyfoot[C]{\thepage}
\renewcommand{\headrulewidth}{0.3pt}
\setlength{\parindent}{1.1em}
\setlength{\parskip}{3pt}
\setcounter{tocdepth}{2}
\numberwithin{equation}{section}
\newtheorem{theorem}{Theorem}[section]
\newtheorem{proposition}[theorem]{Proposition}
\newtheorem{corollary}[theorem]{Corollary}
\newtheorem{lemma}[theorem]{Lemma}
\theoremstyle{definition}
\newtheorem{definition}[theorem]{Definition}
\newtheorem{remark}[theorem]{Remark}
\DeclareMathOperator{\LIS}{LIS}
\DeclareMathOperator{\sgn}{sgn}
\DeclareMathOperator{\E}{\mathbb E}
\DeclareMathOperator{\Prob}{\mathbb P}
\newcommand{\N}{\mathbb N}
\newcommand{\Z}{\mathbb Z}
\newcommand{\Q}{\mathbb Q}
\newcommand{\F}{\mathcal F}
\newcommand{\fall}[2]{(#1)^{\underline{#2}}}
\newcommand{\rise}[2]{(#1)^{\overline{#2}}}
\newcommand{\coef}[2]{[#1^{#2}]}
\newcommand{\oeis}[1]{\href{https://oeis.org/#1}{\texttt{#1}}}
\newcommand{\eps}{\varepsilon}
\lstset{basicstyle=\ttfamily\small,columns=fullflexible,breaklines=true,
 frame=single,rulecolor=\color{ink!30},showstringspaces=false}
% Dated notes added when the manuscript was written into ProveIt (batch 85).
\newenvironment{writenote}[1][]{\par\smallskip\begingroup\small
 \noindent\textbf{[write]}\if\relax\detokenize{#1}\relax\else~\textbf{#1.}\fi\ }%
 {\par\endgroup\smallskip}

\begin{document}
\hypersetup{pageanchor=false}% [write] title page resets the page counter
\begin{titlepage}
\thispagestyle{empty}
\noindent{\color{accent}\large\textsc{Enumerative combinatorics \quad / \quad OEIS research}}
\vspace{0.85in}
\begin{flushleft}
{\color{ink}\Huge\bfseries Finite Fredholm Truncation\par}
\vspace{0.16in}
{\color{ink}\Huge\bfseries for Long Increasing\par}
\vspace{0.10in}
{\color{ink}\Huge\bfseries Subsequences\par}
\vspace{0.4in}
{\Large A proof of fixed-sector holonomicity\par}
\vspace{0.05in}
{\Large and uniform all-orders tail asymptotics\par}
\end{flushleft}
\vspace{0.45in}
\noindent\textbf{Prepared for Vladimir Reshetnikov}\par
\noindent ChatGPT \;\textbar\; October 3, 2026
\vspace{0.45in}
\begin{abstract}
Let $A(N,k)$ count permutations of $N$ letters whose longest increasing
subsequence has length exactly $k$. We prove that, for every fixed integer
$r\ge2$, a globally holonomic array agrees with $A(N,k)$ throughout
$N\le rk$, $k\ge1$. This establishes the fixed-sector existence conjecture
stated by Kauers and Wang in September 2026. The construction uses factorial
moments of the number of shifted Young-diagram rows beyond a threshold.
A rectangle obstruction makes inclusion--exclusion terminate after $r-1$
terms in the sector. The discrete Bessel kernel expresses each fixed moment
as a finite-dimensional proper-hypergeometric sum, so holonomic closure
applies. We determine the exact first failure of each truncation.

We also derive a single-sum formula for the first moment, prove uniform
all-orders asymptotics when $k$ is a positive linear fraction of $N$, and
compute five correction terms for OEIS A269021. The error from replacing
the tail by its first shifted-row moment is either zero or smaller than
every algebraic order. Consequences include rational-ray P-recursiveness,
conditional overshoot estimates, and a mean-subsequence-multiplicity limit.
Exact enumeration and symbolic checks are supplied separately.
\end{abstract}
\vfill
\noindent\fcolorbox{ink!25}{ink!3}{\parbox{0.94\textwidth}{\small
\textbf{Status and scope.} This is an unrefereed conventional mathematical
proof, not a Lean- or Rocq-verified development. The discrete Bessel-kernel
and holonomic-summation theorems are established external inputs. The
specific previously guessed recurrences in the one-third sector are not
certified here; the theorem proved is their underlying existence statement.
Priority beyond the sources checked is not asserted.}}
\end{titlepage}
\hypersetup{pageanchor=true}

\begingroup
\small
\setlength{\parskip}{0pt}
\tableofcontents
\endgroup
\newpage

\section{Problem, provenance, and main results}
\subsection{The arrays and the question}
For $N\ge0$ and $k\ge1$, set
\begin{align}
 A(N,k)&=\#\{\pi\in S_N:\LIS(\pi)=k\},\label{lis:eq:Adef}\\
 T(N,k)&=\#\{\pi\in S_N:\LIS(\pi)\ge k\}.
 \label{lis:eq:Tdef}
\end{align}
Here an increasing subsequence need not occupy adjacent positions. These
are the exact-length array \oeis{A047874} and the upper-tail array
\oeis{A214152}, respectively \cite{oeisA,oeisT}. We use
$A(0,k)=T(0,k)=0$ for $k\ge1$, and both arrays vanish for $k>N$.
They satisfy
\begin{equation}
 A(N,k)=T(N,k)-T(N,k+1).\label{lis:eq:difference}
\end{equation}
The diagonal $T(2n,n)$, with its separately defined empty-permutation value
$a_0=1$, is \oeis{A269021}; $A(2n,n)$, with the same convention at zero,
is \oeis{A267433} \cite{oeis269,oeis267}.

The distinction between a separate recurrence for each fixed $k$ and a
single bivariate recurrence system is essential. A bivariate holonomic
array has uniform polynomial-coefficient recurrences in the index
directions. In particular, there are nonzero operators of the forms
\begin{equation}
 \sum_{i=0}^{d}p_i(N,k)E_N^i,
 \qquad \sum_{i=0}^{e}q_i(N,k)E_k^i,
 \qquad p_i,q_i\in\Q[N,k],\label{lis:eq:pure}
\end{equation}
where $E_NF(N,k)=F(N+1,k)$ and $E_kF(N,k)=F(N,k+1)$.
Orders and degrees do not grow when $N$ and $k$ vary.
We use ``holonomic'' in the discrete holonomic-systems sense, for which
finite sums, products, and definite summation have their usual closure
properties. These properties imply the D-finiteness required here.

Following the regional convention in \cite{KW}, an array is
\emph{D-finite in a region} if it agrees there with a globally D-finite
array. This does not say that a recurrence stencil may leave the region
while its entries are still taken from the original array.

Kauers and Wang proved the half-sector result and the previously
conjectured fourth-order recurrence for A269021 (Conjecture 17 of
\cite{KK}). Their September 2026
paper asks for D-finiteness in every fixed sector $N\le rk$, $r\ge2$
\cite[Sections 1--2]{KW}. We prove this existence statement, rather than
claiming their already established recurrence as new. The main device is
a truncation by the number of shifted rows, not by the number of ordinary
increasing-subsequence witnesses.

\subsection{The main theorem}
For a partition $\lambda=(\lambda_1,\lambda_2,\ldots)$, let $f^\lambda$
be the number of standard Young tableaux of shape $\lambda$. Extend its
parts by zeros and define
\begin{equation}
 Y_k(\lambda)=\#\{i\ge1:\lambda_i-i\ge k-1\},
 \qquad
 C_j(N,k)=\sum_{\lambda\vdash N}(f^\lambda)^2
                   \binom{Y_k(\lambda)}{j}.\label{lis:eq:moments}
\end{equation}
The sum defining $Y_k$ is finite because $k\ge1$.

\begin{theorem}[Finite-moment holonomic extensions]\label{lis:thm:main}
For every fixed $j\ge1$, $C_j(N,k)$ is a holonomic bivariate array.
For fixed $R\ge1$, define
\begin{align}
 T_R(N,k)&=\sum_{j=1}^{R}(-1)^{j+1}C_j(N,k),\label{lis:eq:TR}\\
 A_R(N,k)&=T_R(N,k)-T_R(N,k+1).\label{lis:eq:AR}
\end{align}
Both arrays are globally holonomic and
\begin{equation}
 T_R(N,k)=T(N,k),\qquad A_R(N,k)=A(N,k)
 \quad\text{if}\quad N<(R+1)(k+R).\label{lis:eq:sharpregion}
\end{equation}
Consequently, for every fixed $r\ge2$, the arrays $A$ and $T$ are
D-finite throughout $N\le rk$, $k\ge1$: take $R=r-1$.
\end{theorem}

The word ``globally'' refers to the whole nonnegative $N$, positive $k$
quadrant, not just to the agreement sector. A zero column may be added at
$k=0$ without changing holonomicity. The values of $T_R$ and $A_R$
outside their agreement region are auxiliary values; in particular,
positivity outside the region is not required.

The gain is not merely one additional sector. The same construction
works for every fixed $r$, and its summation dimension depends only on
$r$, not on $N$. At the same time, it does \emph{not} prove global
holonomicity of the original array $A$: the required $R$ is unbounded
as $k/N$ tends to zero.

\begin{theorem}[Uniform tail expansion]\label{lis:thm:asymptotic}
Put $m=N-k$, $\rho=m/k$, and
\begin{equation}
 B(N,k)=\frac{(N!)^2}{(k!)^2(N-k)!}.
 \label{lis:eq:base}
\end{equation}
For every compact interval $I\subset(0,\infty)$ and integer $M\ge0$,
uniformly over integer pairs with $k\to\infty$ and $\rho\in I$,
\begin{equation}
 \frac{T(N,k)}{B(N,k)}
 =e^{-2\rho}\left(\sum_{d=0}^{M}\frac{c_d(\rho)}{k^d}
                    +O_{I,M}(k^{-M-1})\right).\label{lis:eq:allorders}
\end{equation}
The coefficients are explicitly computable polynomials, $c_0=1$,
$\deg c_d\le2d$, and
\begin{align}
 c_1(\rho)&=2\rho-\rho^2,\label{lis:eq:c1}\\
 c_2(\rho)&=\frac{\rho(3\rho^3-20\rho^2+42\rho+2)}6,\label{lis:eq:c2}\\
 c_3(\rho)&=-\frac{\rho(\rho^5-14\rho^4+73\rho^3-158\rho^2
                                  +26\rho+30)}6.\label{lis:eq:c3}
\end{align}
An all-orders coefficient rule is given in Section~\ref{lis:sec:coefficients}.
\end{theorem}

These are Poincar\'e asymptotic expansions with a remainder for every
fixed truncation order, not a claim of convergence of the infinite
series. Section~\ref{lis:sec:firsterror} identifies a distinct,
beyond-all-algebraic-orders contribution from diagrams with a second
long shifted row. We do not claim a full expansion of that exponentially
small contribution.

\subsection{Relationship to earlier work and the repository}
The leading A269021 asymptotic $16^n(n-1)!/(\pi e^2)$ is already recorded
in the OEIS and attributed to V\'aclav Kotesovec in 2016 \cite{oeis269}.
The same leading equivalent for A267433 has a public proof by David
Moews from 2023 \cite{Moews}. The contribution here is the uniform
all-orders development, together with the fixed-sector theorem and
its sharp truncation boundary, not the discovery of that leading
constant.

The ProveIt repository was inspected at revision
\texttt{6bf7f30d0352f7596e70928b3d4f304914075907} \cite{repo}.
A targeted identifier search did not surface an A269021 report. That is
a bounded search observation, not an exhaustive assertion about every
artifact or an external priority search. This article uses no unverified
mathematical assertion from unrelated repository projects. It is designed
as a self-contained research package suitable for later integration.

The new deductions below rely explicitly on three standard inputs:
the Robinson--Schensted enumeration, the discrete Bessel correlation
formula, and holonomic summation. The remaining reductions, cutoff
analysis, finite-sum formulas, and asymptotic error estimates are proved
in the article. The package includes no claim of minimal recurrence order.

\begin{writenote}[Added 3 October 2026, batch~85]
\emph{Provenance.} This report is a single manuscript, printed in full:
manuscript~08 of batch~85 of ProveIt's incoming-reports intake, delivered
as the archive \texttt{LIS\_\allowbreak fixed\_\allowbreak
sector\_\allowbreak research\_\allowbreak package.zip} (wrapper directory
\texttt{LIS\_\allowbreak fixed\_\allowbreak sector\_\allowbreak
research/}) in commit \texttt{9d6968c8a}, and placed in commit
\texttt{ddf8df5d5} in the research-report collection, among the
OEIS-sequence asymptotics reports, as the report
\texttt{a047874-\allowbreak long-\allowbreak increasing-\allowbreak
subsequences}, named after the array whose sector question it answers.
Its pin is the revision \texttt{6bf7f30d0} quoted above. The title page
and the PDF author field read ``ChatGPT'': the package names ChatGPT as
the tool that wrote it and names no human author. Every label now carries
the prefix \texttt{lis:}; no statement, proof or number of the manuscript
was changed.

\emph{What Kauers and Wang state.} The arXiv abstract of \cite{KW} does
not mention the sector question, so on 3~October 2026 the intake read the
HTML text of arXiv:2609.02220v1 (licensed CC~BY~4.0). Its Section~1
defines ``D-finite inside a region'' as the existence of a D-finite array
coinciding with the original one there (the convention used above),
proves two bivariate recurrences valid for $n/2\le k$ (their Theorem~1,
for their $a_{n,k}$, which is $A(N,k)$ here) and derives
Conjecture~17 of \cite{KK} from them. At the end of their Section~2,
after displaying guessed recurrences of higher order and degree for
$n/3<k<n/2$, they write: ``We suspect that for every $r=2,3,\dots$, the
sequence $a_{n,k}$ is D-finite in the range $k\geq n/r$''; they add that
this would not give D-finiteness for all $k\ge0$, and hope for a uniform
approach to all such slices. So the statement is an informal conjecture,
not a numbered one; ``the fixed-sector existence conjecture'' is this
manuscript's name for it. Its content is exactly the sector assertion of
Theorem~\ref{lis:thm:main} ($k\ge n/r$ is $N\le rk$). For $r=2$ the
D-finiteness is already Kauers and Wang's Theorem~1, with explicit
recurrences; there Theorem~\ref{lis:thm:main} is only a second route to
existence. The new content is every $r\ge3$ and the affine cutoff
\eqref{lis:eq:sharpregion}. Kauers and Wang's guessed $r=3$ operators and
all minimal orders remain uncertified (Section~4.3).

\emph{Review status.} Nothing here is refereed or formalized. The
holonomic-closure step of Section~3.4 is the step most in need of
independent review; see the note there. The intake checked the finite
identities independently: a separate hook-length program reproduced
\eqref{lis:eq:C1sum} for all $N\le14$, $1\le k\le N+1$, the truncation
identity \eqref{lis:eq:sharpregion} and the first failure
\eqref{lis:eq:firstfailure} for $R=1,2,3$ in that range, and the
diagonal values of the table in Section~9.2 for $n\le6$.

\emph{Repository.} The statement above that a targeted search ``did not
surface an A269021 report'' is true at the pin and was confirmed at
placement: no other file of ProveIt treats A047874, A214152, A269021 or
A267433, and no Lean or Rocq development treats Young tableaux or
longest increasing subsequences. Other conjectures of \cite{KK} are
treated in the collection reports
\texttt{enumerative-\allowbreak combinatorics/\allowbreak
a181280-\allowbreak binary-\allowbreak matrix-\allowbreak formula}
(Conjecture~20), \texttt{enumerative-\allowbreak combinatorics/\allowbreak
a195806-\allowbreak hexagonal-\allowbreak lattice} (Conjecture~11) and
Part~IV of \texttt{a215561-\allowbreak fixed-\allowbreak
composition-\allowbreak excursions} (Conjecture~15); Conjecture~17 is
Kauers and Wang's, not this report's. The batch-85 sibling
\texttt{a181199-\allowbreak shifted-\allowbreak rectangles} treats
fixed-height shifted rectangles. The rectangle formula
\eqref{lis:eq:rectanglehook} is the classical hook-length count, printed
also as \texttt{t2:eq:tableaux} in the Transseries volume
\texttt{Combinatorial\_\allowbreak Transseries\_\allowbreak Inverses}.

\emph{Notation.} No symbol was renamed. Some letters carry two meanings:
\begin{center}\footnotesize
\begin{tabular}{@{}>{\raggedright\arraybackslash}p{0.17\textwidth}>{\raggedright\arraybackslash}p{0.36\textwidth}>{\raggedright\arraybackslash}p{0.39\textwidth}@{}}
\toprule
Symbol & Meaning here & Other meaning, or tempting false reading\\
\midrule
$A(N,k)$, $T(N,k)$ & exact-length and tail arrays (A047874, A214152) &
Kauers and Wang write $a_{n,k}$ and $b_{n,k}$\\
$a_n$, $b_n$ (Section~9) & $a_n=T(2n,n)$ (A269021), $b_n=A(2n,n)$
(A267433) & not Kauers and Wang's arrays: their $b_{2n,n}$ is $a_n$ here\\
$R$ & truncation rank & Kauers and Wang's name for a region\\
$j$ & moment order of $C_j$ & summation index in
\eqref{lis:eq:A269single}--\eqref{lis:eq:A267single} and Section~\ref{lis:sec:coefficients}\\
$m$ & summation index in Lemma~\ref{lis:lem:collapse} and
Theorem~\ref{lis:thm:single} & the deficit $N-k$ in
Theorem~\ref{lis:thm:asymptotic} and Section~\ref{lis:sec:coefficients}\\
$q$ & $N/(k+2)^2$ in Proposition~\ref{lis:prop:errorbound} & the
summation variables $q_a$ of Theorem~\ref{lis:thm:multisum}; the
denominator in Corollary~\ref{lis:cor:rays}\\
$C_1$ & first shifted-row moment & not the mean witness count
(Section~5.2); not the constants $C_I$\\
$L_\ell$, $L_n$ & logarithmic coefficients \eqref{lis:eq:Lell} &
$L_n$ is the leading term of $a_n$ in Section~9.3\\
\bottomrule
\end{tabular}
\end{center}
\end{writenote}

\section{Shifted rows and exact inclusion--exclusion}
\subsection{From permutations to tableaux}
The Robinson--Schensted correspondence gives a bijection between
permutations and pairs of standard Young tableaux of the same shape,
with the longest increasing subsequence equal to the first row length
\cite{Schensted}. Hence
\begin{equation}
 T(N,k)=\sum_{\substack{\lambda\vdash N\\\lambda_1\ge k}}
                    (f^\lambda)^2,
 \qquad \sum_{\lambda\vdash N}(f^\lambda)^2=N!.
 \label{lis:eq:RSK}
\end{equation}
We use the hook-length formula to compute $f^\lambda$ in the verifier;
the structural proof only needs \eqref{lis:eq:RSK} and the standard rectangle
dimension formula when describing the first discrepancy.

The shifted row coordinates $\lambda_i-i$ are strictly decreasing.
Therefore $Y_k\ge1$ is precisely $\lambda_1\ge k$. More generally:

\begin{lemma}[Rectangle obstruction]\label{lis:lem:rectangle}
For $j,k\ge1$,
\begin{equation}
 Y_k(\lambda)\ge j
 \quad\Longleftrightarrow\quad
 \lambda_j\ge k+j-1.
 \label{lis:eq:rectangleequiv}
\end{equation}
Thus $Y_k\ge j$ forces $|\lambda|\ge j(k+j-1)$. At equality, the only
possible partition is the rectangle $(k+j-1)^j$.
\end{lemma}
\begin{proof}
Strict decrease of the shifted coordinates means that at least $j$
coordinates meet the threshold exactly when the $j$th one does.
The first $j$ rows then each have at least $k+j-1$ boxes. Equality in
the total size forces these rows to have exactly that length and
leaves no further boxes. Conversely, that rectangle meets the threshold.
\end{proof}

This is the source of the correction $j(j-1)$ in all cutoff formulas.
Counting ordinary rows of length at least $k$ would give the wrong
threshold and the wrong moment process.

\subsection{The residual identity and its sign}
For integers $y\ge1$ and $R\ge1$, Pascal's identity gives
\begin{equation}
 \sum_{j=1}^{R}(-1)^{j+1}\binom yj
 =1-(-1)^R\binom{y-1}{R}.
 \label{lis:eq:binomialresidual}
\end{equation}
The identity remains useful when $R\ge y$, using zero for an
out-of-range binomial coefficient. For $y=0$, the sum is zero.
Weighting \eqref{lis:eq:binomialresidual} by $(f^\lambda)^2$ proves the
following exact formula.

\begin{proposition}[Bonferroni residual]\label{lis:prop:residual}
For all $N\ge0$, $k,R\ge1$,
\begin{equation}
 T_R(N,k)-T(N,k)
 =(-1)^{R+1}
   \sum_{\substack{\lambda\vdash N\\Y_k(\lambda)\ge1}}
       (f^\lambda)^2\binom{Y_k(\lambda)-1}{R}.
 \label{lis:eq:residual}
\end{equation}
In particular, odd truncations are upper bounds and even truncations
are lower bounds for $T$. The residual vanishes if
$N<(R+1)(k+R)$.
\end{proposition}
\begin{proof}
Equation~\eqref{lis:eq:RSK} identifies the contribution of the constant $1$
on the right side of \eqref{lis:eq:binomialresidual} with $T$. A nonzero
residual requires $Y_k\ge R+1$, and Lemma~\ref{lis:lem:rectangle} excludes
this below the stated size. Every weight in the residual sum is
nonnegative, proving its sign.
\end{proof}

The sign statement belongs to the tail truncations. A difference of two
Bonferroni bounds need not give a bound of the same sign for the
exact-length array away from the special first-discrepancy point.

\begin{proposition}[The first failure is exact and nonzero]\label{lis:prop:firstfailure}
Let $N_0=(R+1)(k+R)$ and $\mu=(k+R)^{R+1}$. Then
\begin{equation}
 T_R(N_0,k)-T(N_0,k)
 =A_R(N_0,k)-A(N_0,k)
 =(-1)^{R+1}(f^\mu)^2.\label{lis:eq:firstfailure}
\end{equation}
For a rectangle $b^a$,
\begin{equation}
 f^{b^a}=(ab)!\,\frac{\prod_{i=0}^{a-1}i!}
                           {\prod_{i=0}^{a-1}(b+i)!}.
 \label{lis:eq:rectanglehook}
\end{equation}
Thus \eqref{lis:eq:sharpregion} is the sharp universal cutoff for this
particular truncation.
\end{proposition}
\begin{proof}
At $N_0$, the rectangle in Lemma~\ref{lis:lem:rectangle} is the only shape
with $Y_k\ge R+1$. It has $Y_k=R+1$, so the residual binomial coefficient
is one. At threshold $k+1$ the first possible omitted rectangle is
larger by $R+1$, so $T_R(N_0,k+1)=T(N_0,k+1)$. Taking differences
proves the assertion for $A_R$. Formula~\eqref{lis:eq:rectanglehook}
is the hook-length formula specialized to a rectangle.
\end{proof}

\begin{center}
\begin{tabular}{@{}cll@{}}
\toprule
$R$ & Exact tail expression & Region of exactness\\
\midrule
1 & $C_1$ & $N<2k+2$\\
2 & $C_1-C_2$ & $N<3k+6$\\
3 & $C_1-C_2+C_3$ & $N<4k+12$\\
4 & $C_1-C_2+C_3-C_4$ & $N<5k+20$\\
\bottomrule
\end{tabular}
\end{center}
In the first line, the error at $N=2k+2$ is the square of the
$(k+1)$st Catalan number, since $f^{(k+1,k+1)}$ equals that number.
The affine cutoff is stronger than the homogeneous sector required
by the conjecture.

\section{The Bessel kernel and a fixed-dimensional factorial sum}
\subsection{The established determinantal input}
For $t>0$, poissonized Plancherel measure is
\begin{equation}
 \Prob_t(\lambda)=e^{-t}t^{|\lambda|}
                 \frac{(f^\lambda)^2}{(|\lambda|!)^2}.
 \label{lis:eq:poisson}
\end{equation}
It is a probability measure by \eqref{lis:eq:RSK}. Let
$D(\lambda)=\{\lambda_i-i:i\ge1\}\subset\Z$.
The correlation theorem of Borodin, Okounkov, and Olshanski
\cite[Theorem 2 and Proposition 2.9]{BOO} states that
\begin{equation}
 \Prob_t(\{x_1,\ldots,x_j\}\subset D(\lambda))
       =\det[K_t(x_a,x_b)]_{a,b=1}^{j},\label{lis:eq:kernelcorr}
\end{equation}
for distinct $x_a$, where
\begin{equation}
 K_t(x,y)=\sum_{\ell=1}^{\infty}
       J_{x+\ell}(2\sqrt t)J_{y+\ell}(2\sqrt t).
 \label{lis:eq:kernel}
\end{equation}
Only nonnegative $x,y$ are used here. The Bessel series needed below is
\begin{equation}
 J_d(2\sqrt t)=t^{d/2}\sum_{s\ge0}
                  \frac{(-1)^st^s}{s!(d+s)!},\qquad d\ge0.
 \label{lis:eq:Bessel}
\end{equation}
The normalization in \eqref{lis:eq:poisson} contains $(N!)^2$, not $N!$.
This factor must be restored after coefficient extraction.

Summing correlations over increasing $j$-tuples in the threshold set
$\{k-1,k,\ldots\}$ gives its factorial moment:
\begin{equation}
 e^{-t}\sum_{N\ge0}\frac{C_j(N,k)}{(N!)^2}t^N
 =\sum_{k-1\le x_1<\cdots<x_j}
                     \det[K_t(x_a,x_b)].\label{lis:eq:momentgf}
\end{equation}
There is no $1/j!$ here: the $x$-tuples are strictly ordered and hence
already represent unordered subsets once each.

\subsection{Cauchy--Binet removes the kernel summation}
Write $x_a=k-1+u_a$ and $\ell_b=1+v_b$. Applying Cauchy--Binet to
\eqref{lis:eq:kernel} yields
\begin{equation}
 e^{-t}\sum_{N\ge0}\frac{C_j(N,k)}{(N!)^2}t^N
 =\sum_{\substack{0\le u_1<\cdots<u_j\\0\le v_1<\cdots<v_j}}
       \left(\det[J_{k+u_a+v_b}(2\sqrt t)]_{a,b=1}^{j}\right)^2.
 \label{lis:eq:squaredminors}
\end{equation}
This also explains the Fredholm terminology: these are the elementary
symmetric, or Fredholm-expansion, coefficients of the threshold kernel.
We truncate the expansion after a fixed number of these coefficients;
we do \emph{not} assert that the infinite kernel has finite operator rank.

All operations in \eqref{lis:eq:squaredminors} are coefficientwise finite.
Indeed, factoring the powers of $t^{1/2}$ from rows and columns shows
that a squared minor has lowest possible degree
\begin{equation}
 b=jk+\sum_a u_a+\sum_b v_b.
 \label{lis:eq:basedegree}
\end{equation}
Since $u_a,v_a\ge a-1$, this is at least $j(k+j-1)$.
For a fixed degree $N$, only finitely many $u,v$ have $b\le N$.
Although an individual Bessel factor may involve a half-integral power,
each squared minor has integral powers. Equivalently, work first with
$z=\sqrt t$ and retain its even powers.

For analytic justification, finite Cauchy--Binet is applied before taking
larger index cutoffs. The Bessel series converges absolutely on compact
sets and its orders have factorial decay. Thus the finite cutoffs converge
locally uniformly. Alternatively, every coefficient stabilizes once the
cutoffs exceed its finite degree budget, which gives the same formal
identity. No unproved interchange of an asymptotic expansion with an
infinite determinant is being used.

\subsection{The explicit multisum}
For permutations $\sigma,\tau\in S_j$ and nonnegative integer vectors
$s=(s_1,\ldots,s_j)$ and $q=(q_1,\ldots,q_j)$, put
\begin{equation}
 h=N-jk-\sum_a u_a-\sum_b v_b-\sum_a(s_a+q_a).
 \label{lis:eq:h}
\end{equation}

\begin{theorem}[Proper-hypergeometric formula]\label{lis:thm:multisum}
For fixed $j\ge1$, the following finite sum holds for all $N\ge0$, $k\ge1$:
\begin{align}
 \frac{C_j(N,k)}{(N!)^2}
 &=\sum_{\sigma,\tau\in S_j}\sgn(\sigma)\sgn(\tau)
   \sum_{\substack{0\le u_1<\cdots<u_j,\ 0\le v_1<\cdots<v_j\\
                    s_a,q_a\ge0,\ h\ge0}}
      \frac{(-1)^{\sum_a(s_a+q_a)}}{h!}\notag\\[-2pt]
 &\hspace{1.0cm}\cdot
   \prod_{a=1}^{j}\frac{1}{s_a!q_a!
          (k+u_a+v_{\sigma(a)}+s_a)!
          (k+u_a+v_{\tau(a)}+q_a)!}.
 \label{lis:eq:multisum}
\end{align}
For each pair $(\sigma,\tau)$ there are $4j$ summation variables,
independently of $N$ and $k$.
\end{theorem}
\begin{proof}
Expand the two determinants in \eqref{lis:eq:squaredminors}, and then use
\eqref{lis:eq:Bessel} in each of their $2j$ factors. The total initial
power is \eqref{lis:eq:basedegree}, because
$\sum_a v_{\sigma(a)}=\sum_a v_{\tau(a)}=\sum_bv_b$.
Multiplication by $e^t$ inverts the $e^{-t}$ in \eqref{lis:eq:squaredminors}.
Its coefficient supplying the remaining degree is $1/h!$. This gives
\eqref{lis:eq:multisum}. Every summation index is bounded once $h\ge0$
and \eqref{lis:eq:h} are imposed, so the calculation is an identity of
finite coefficient sums.
\end{proof}

\subsection{Why holonomic summation applies, including the boundaries}
We use the standard proper-hypergeometric summation theorem of
Wilf--Zeilberger \cite[Section 2, Theorem 1 and Corollary A]{WZ}, together
with definite-summation closure for holonomic systems
\cite[Theorem 2 and Section 4]{Koutschan}. In the form needed here:
\begin{quote}
A fixed-dimensional finite sum of proper-hypergeometric terms, with
factorials of integer-affine forms and integer-linear summation boundaries,
is holonomic in its remaining integer parameters. Finite sums and
products of the resulting holonomic arrays remain holonomic.
\end{quote}
A proper term has factorial products or quotients, polynomial factors,
and constant-to-integer powers; mere rationality of shift quotients is
not the only hypothesis being invoked.

To apply this statement concretely, remove strict inequalities by the
gap coordinates
\begin{equation}
 u_a=a-1+\sum_{b=1}^{a}\alpha_b,
 \qquad v_a=a-1+\sum_{b=1}^{a}\beta_b,
 \qquad \alpha_b,\beta_b\ge0.
 \label{lis:eq:gaps}
\end{equation}
Then
\begin{equation}
 h=N-jk-j(j-1)-\sum_{b=1}^{j}(j-b+1)(\alpha_b+\beta_b)
                        -\sum_a(s_a+q_a).
 \label{lis:eq:hgap}
\end{equation}
For fixed $j$, every factorial argument in \eqref{lis:eq:multisum} is an
integer-affine expression in
\[
 N,\ k,\ \alpha_1,\ldots,\alpha_j,\beta_1,\ldots,\beta_j,
       s_1,\ldots,s_j,q_1,\ldots,q_j.
\]
In particular, $jk$ is linear because $j$ is fixed.

One can regard the inner sum as a sum over the cube $\{0,\ldots,N\}^{4j}$,
setting $1/h!=0$ for negative integer $h$. No nonzero term is lost,
since every coefficient in the budget \eqref{lis:eq:hgap} is positive.
This reciprocal-factorial zero convention is the usual terminating
hypergeometric convention. The reciprocal-gamma recurrence extends
across its integer zeros. Restriction to the nonnegative cube gives
ordinary holonomic summation boundaries; alternatively, its indicator
has a rational generating function. Boundary terms produced by shifting
$N$ or $k$ are included by definite-summation closure and can be
homogenized. We do not simply discard such boundary terms.

Thus the closure theorem applies simultaneously with $N,k$ left as
parameters. Multiplying by $(N!)^2$ preserves holonomicity, as do the
finitely many $\sigma,\tau$ terms. This proves the first assertion of
Theorem~\ref{lis:thm:main}. It is stronger than proving P-recursiveness
separately after specializing $k$.

\begin{writenote}[Added 3 October 2026, batch~85: needs independent review]
This subsection carries the whole existence claim, and the intake did not
verify it; it should be reviewed independently before the sector theorem
is relied on (the package's own audit also says that independent review
``remains appropriate''). A reviewer should check at least: that the
summand of \eqref{lis:eq:multisum}, in the gap coordinates
\eqref{lis:eq:gaps}, is holonomic as a function of all $4j+2$ integer
variables, $N$ and $k$ included, under the stated conventions (the
reciprocal factorial $1/h!$ vanishing for $h<0$, and the restriction to
$\alpha_b,\beta_b,s_a,q_a\ge0$, which is imposed by a cube rather than by
reciprocal factorials for $\alpha_b,\beta_b$); that the closure theorem
cited from \cite{WZ,Koutschan} applies with the $N$-dependent support
cut out by $h\ge0$; and that the boundary terms of creative telescoping,
said above to be ``included by definite-summation closure and [able to]
be homogenized'', are in fact accounted for. The executable checks test
the identity of Theorem~\ref{lis:thm:multisum} (through the 140
Bessel-minor comparisons) and the finite truncation identities on
recorded ranges; they test no holonomicity. Sections~5--9 do not depend
on this step: the single sums, the uniform expansion and the A269021 and
A267433 corrections use Theorem~\ref{lis:thm:single},
Proposition~\ref{lis:prop:residual} and
Proposition~\ref{lis:prop:errorbound}, not holonomic closure. The
rational-ray Corollary~\ref{lis:cor:rays} of Section~4 does depend on it.
\end{writenote}

\begin{proof}[Completion of Theorem~\ref{lis:thm:main}]
Finite linear combinations and shifts give holonomicity of $T_R$ and
$A_R$. Proposition~\ref{lis:prop:residual} gives agreement of $T_R$ with $T$
below $(R+1)(k+R)$. The same inequality also implies agreement at $k+1$,
so \eqref{lis:eq:difference} gives agreement for $A_R$.
For $R=r-1$ and $N\le rk$,
\[
 N\le rk<r(k+r-1)=(R+1)(k+R),
\]
which proves the sector assertion.
\end{proof}

\section{Consequences for recurrence theory}
\subsection{Rational rays and periodic rounding}
\begin{corollary}[Every positive rational ray]\label{lis:cor:rays}
Fix integers $a\ge b\ge1$ and offsets $a_0,b_0$.
The sequences
\[
 A(an+a_0,bn+b_0),\qquad T(an+a_0,bn+b_0)
\]
are eventually P-recursive, and any finite completion at small indices
is P-recursive. The same conclusion holds for
$A(N,\lfloor pN/q\rfloor+c)$ and $T(N,\lfloor pN/q\rfloor+c)$,
where $0<p\le q$ and $c$ are fixed integers.
\end{corollary}
\begin{proof}
Choose fixed $R$ with $a<(R+1)b$. The ray lies in the exactness region
for all sufficiently large $n$. Substitution in \eqref{lis:eq:multisum}
leaves a fixed-dimensional proper-hypergeometric sum with integer-affine
factorial arguments in $n$ and the summation indices. Hence its sequence
is P-recursive by the same summation theorem. A finite modification
preserves P-recursiveness. For the rounded ray, split $N$ into its $q$
residue classes. Each subsequence is a linear ray, and interlacing finitely
many P-recursive sequences preserves P-recursiveness: their generating
functions combine as $\sum_{s=0}^{q-1}z^sF_s(z^q)$.
\end{proof}

No assertion about irrational slopes follows from this argument.
Nor does it impose an effective small order on a recurrence along a ray.
The dimension is fixed but can still be computationally expensive.

\subsection{An abstract transfer principle}
The argument isolates a reusable fact. Suppose a nonnegative integer
statistic $Y$ has weighted factorial moments $C_j(N,k)$ that are
holonomic for every fixed $j$, and suppose $Y\ge j$ is impossible
below a known size $w_j(k)$. Then the weighted event count for $Y\ge1$
is represented by a holonomic array on $N<w_{R+1}(k)$, for each fixed
$R$. The proof is exactly \eqref{lis:eq:binomialresidual}.

The nontrivial application here is that the determinantal correlations
make the moments fixed-dimensional even while $k$ grows, and that the
Young-diagram obstruction provides the explicit function
$w_j(k)=j(k+j-1)$. A growing-size determinant by itself would not
justify bivariate holonomicity.

\subsection{What has and has not been certified}
The construction proves the existence of the required uniform
recurrences and supplies a factorial-sum input from which creative
telescoping can, in principle, compute them. It does not show that any
particular previously guessed operator belongs to the annihilator.
To certify such an operator, one must produce telescoping certificates
or compare a sufficiently complete holonomic system with controlled
initial and singular data. Checking many values alone is not that proof.

A useful next computation would use $C_1-C_2$ rather than the entire
original array. This has a fixed eight-variable inner sum for $C_2$,
with four signed determinant-expansion terms. Cancellation, determinant
structure, or symmetry may reduce this substantially before elimination.
The present existence proof does not depend on successfully performing
that large elimination.

\section{A single sum for the first moment}
\label{lis:sec:single}
\subsection{Collapsing the trace}
At $j=1$, \eqref{lis:eq:squaredminors} simplifies to
\begin{equation}
 e^{-t}\sum_{N\ge0}\frac{C_1(N,k)}{(N!)^2}t^N
 =\sum_{d\ge k}(d-k+1)J_d(2\sqrt t)^2.
 \label{lis:eq:trace}
\end{equation}
Indeed, exactly $d-k+1$ pairs $(u,v)\in\Z_{\ge0}^2$ satisfy
$u+v=d-k$.

\begin{lemma}[Two elementary coefficient identities]\label{lis:lem:collapse}
For integers $m\ge d\ge0$,
\begin{equation}
 \coef{t}{m}J_d(2\sqrt t)^2
 =\frac{(-1)^{m-d}}{(m!)^2}\binom{2m}{m-d}.
 \label{lis:eq:besselsquare}
\end{equation}
For $m\ge k\ge1$,
\begin{equation}
 \sum_{d=k}^{m}(d-k+1)(-1)^{m-d}\binom{2m}{m-d}
 =(-1)^{m-k}\binom{2m-2}{m-k}.
 \label{lis:eq:alternatingbinom}
\end{equation}
\end{lemma}
\begin{proof}
Squaring \eqref{lis:eq:Bessel} gives a sum over $s+q=m-d$.
After multiplication by $(m!)^2$, its unsigned summand is
$\binom ms\binom m{m-d-s}$. Vandermonde's convolution gives
\eqref{lis:eq:besselsquare}. For \eqref{lis:eq:alternatingbinom}, set
$L=m-k$ and $a=m-d$. The left side is
\[
 \sum_{a=0}^{L}(L-a+1)(-1)^a\binom{2m}{a}
 =\coef{x}{L}\frac{(1-x)^{2m}}{(1-x)^2},
\]
which is the right side.
\end{proof}

\begin{theorem}[First-moment single sum]\label{lis:thm:single}
For every $N\ge0$ and $k\ge1$,
\begin{equation}
 C_1(N,k)=(N!)^2\sum_{m=k}^{N}
       \frac{(-1)^{m-k}\binom{2m-2}{m-k}}{(m!)^2(N-m)!}.
 \label{lis:eq:C1sum}
\end{equation}
It equals $T(N,k)$ whenever $N\le2k+1$. In the same region,
\begin{equation}
 A(N,k)=(N!)^2\sum_{m=k}^{N}
       \frac{(-1)^{m-k}\binom{2m-1}{m-k}}{(m!)^2(N-m)!}.
 \label{lis:eq:Asum}
\end{equation}
\end{theorem}
\begin{proof}
Insert Lemma~\ref{lis:lem:collapse} in \eqref{lis:eq:trace} and multiply by
$e^t$. Coefficient extraction gives \eqref{lis:eq:C1sum}.
Proposition~\ref{lis:prop:residual} gives the agreement region.
Subtract the $k+1$ sum from the $k$ sum; their opposite signs and
Pascal's identity give \eqref{lis:eq:Asum}.
\end{proof}

In particular, for $n\ge1$,
\begin{align}
 T(2n,n)&=((2n)!)^2\sum_{j=0}^{n}
  \frac{(-1)^j\binom{2n+2j-2}{j}}{((n+j)!)^2(n-j)!},
  \label{lis:eq:A269single}\\
 A(2n,n)&=((2n)!)^2\sum_{j=0}^{n}
  \frac{(-1)^j\binom{2n+2j-1}{j}}{((n+j)!)^2(n-j)!}.
  \label{lis:eq:A267single}
\end{align}
These formulas are suitable for exact rational arithmetic, with an
integer result. They provide an independent computational route to the
two OEIS diagonals without enumerating all partitions of $2n$.

\subsection{A normalization with a combinatorial meaning}
The quantity $B(N,k)$ in \eqref{lis:eq:base} can also be written
\begin{equation}
 B(N,k)=\binom Nk^2(N-k)!.
 \label{lis:eq:Bwitness}
\end{equation}
It counts pairs consisting of a permutation and a distinguished
increasing subsequence of length $k$: choose its positions and its
values, place those values in increasing order, and permute the
remaining values arbitrarily.

Thus $T/B$ is the reciprocal of the mean number of such witnesses
conditional on their existence. The auxiliary moment $C_1$, however,
counts shifted rows, not witnesses. Confusing these two first moments
would remove the factor $e^{-2\rho}$ from the asymptotic answer.

\section{Why higher shifted-row moments are beyond all orders}
\label{lis:sec:firsterror}
The exact first-moment formula remains useful beyond its exact tail
region. The next estimate controls the discrepancy without a contour
integral or a saddle-point argument.

\begin{proposition}[Explicit first-moment error bound]\label{lis:prop:errorbound}
If $N<2k+2$, then $C_1(N,k)=T(N,k)$.
If $N\ge2k+2$ and $q=N/(k+2)^2<1$, then
\begin{equation}
 0\le\frac{C_1(N,k)-T(N,k)}{B(N,k)}
 \le \frac{N}{k}(1-q)^{-2}
       \frac{(k!)^2(N-k)!}{((k+1)!)^4(N-2k-2)!}.
 \label{lis:eq:errorbound}
\end{equation}
For $\rho=(N-k)/k$ in any fixed compact $I\subset(0,\infty)$,
there are constants $C_I,c_I>0$ such that, for sufficiently large $N$,
\begin{equation}
 0\le\frac{C_1(N,k)-T(N,k)}{B(N,k)}
 \le C_I\exp(-c_I N\log N).
 \label{lis:eq:superalgebraic}
\end{equation}
\end{proposition}
\begin{proof}
The $R=1$ residual is a sum of $(f^\lambda)^2(Y_k-1)$ over $Y_k\ge1$.
It is nonzero only if $Y_k\ge2$, equivalently $\lambda_2\ge k+1$.
Also $Y_k\le N/k$, by Lemma~\ref{lis:lem:rectangle}. Therefore
\begin{equation}
 0\le C_1-T\le\frac Nk
       \sum_{\substack{\lambda\vdash N\\\lambda_2\ge k+1}}
                      (f^\lambda)^2.\label{lis:eq:tworowbound}
\end{equation}
Write $\lambda=(\ell_1,\ell_2,\nu)$ and $v=N-\ell_1-\ell_2$.
Choose the labels for the first two rows and for the remaining diagram,
and ignore all column restrictions involving the first two rows.
The rows have forced increasing order after their labels are chosen,
so
\[
 f^{(\ell_1,\ell_2,\nu)}
 \le \frac{N!}{\ell_1!\ell_2!v!}\,f^\nu.
\]
Using $\sum_{\nu\vdash v}(f^\nu)^2=v!$ and discarding restrictions on
$\nu$ gives
\begin{equation}
 \sum_{\lambda_2\ge k+1}(f^\lambda)^2
 \le (N!)^2\sum_{\substack{\ell_1,\ell_2\ge k+1\\\ell_1+\ell_2\le N}}
        \frac{1}{(\ell_1!)^2(\ell_2!)^2(N-\ell_1-\ell_2)!}.
 \label{lis:eq:doublemajorant}
\end{equation}
We even dropped $\ell_1\ge\ell_2$, which can only enlarge the sum.
Increasing either $\ell_i$ by one multiplies its summand by at most
$N/(k+2)^2=q$. The double sum is consequently bounded by its
$\ell_1=\ell_2=k+1$ term times $(1-q)^{-2}$. Substituting this bound
in \eqref{lis:eq:tworowbound} and dividing by $B$ proves
\eqref{lis:eq:errorbound}.

For the uniform conclusion, write $k\ge\alpha_0N$ with
$\alpha_0>0$ depending only on $I$. Then $q=O_I(N^{-1})$, and
\begin{equation}
 \frac{(k!)^2(N-k)!}{((k+1)!)^4(N-2k-2)!}
 \le \frac{N^{k+2}}{(k!)^2(k+1)^4}
 \le \frac{N^2}{(k+1)^4}
                 \left(\frac{e^2N}{k^2}\right)^k.
 \label{lis:eq:factorialmajorant}
\end{equation}
Here $k!\ge(k/e)^k$ follows by comparing $\sum_{i=1}^k\log i$
with the integral of $\log x$. The logarithm of the right side of
\eqref{lis:eq:factorialmajorant} is at most
$-\alpha_0N\log N+O_I(N)$. This implies
\eqref{lis:eq:superalgebraic}; in the other region the error is exactly zero.
\end{proof}

The estimate separates two logically different issues. The holonomic
sector theorem requires several moments to obtain an \emph{exact}
answer. The all-orders algebraic asymptotic expansion requires only the
first moment, because every omitted contribution is negligible compared
with every power of $k^{-1}$. This explains why the same asymptotic
coefficient family works on both sides of $N=2k$.

\section{Uniform all-orders expansion and its coefficient calculus}
\label{lis:sec:coefficients}
\subsection{The normalized finite product}
Use $m=N-k$ for the deficit and replace the summation index in
\eqref{lis:eq:C1sum} by $k+j$. Division by $B$ gives
\begin{equation}
 \frac{C_1(N,k)}{B(N,k)}
 =\sum_{j=0}^{m}(-1)^j\binom{2k+2j-2}{j}
                         \frac{\fall{m}{j}}{\bigl(\rise{k+1}{j}\bigr)^2}.
 \label{lis:eq:normalizedsum}
\end{equation}
Writing $m=\rho k$, this is
\begin{equation}
 \frac{C_1}{B}=\sum_{j=0}^{m}\frac{(-2\rho)^j}{j!}
                       Q_{k,j}(\rho),\label{lis:eq:Qsum}
\end{equation}
where
\begin{equation}
 Q_{k,j}(\rho)=
 \frac{\displaystyle\prod_{a=0}^{j-1}
    \left(1+\frac{2j-2-a}{2k}\right)
    \left(1-\frac{a}{\rho k}\right)}
      {\displaystyle\prod_{a=1}^{j}\left(1+\frac ak\right)^2}.
 \label{lis:eq:Qproduct}
\end{equation}
Empty products are one. For each fixed $j$, $Q_{k,j}\to1$.
The alternating exponential series in \eqref{lis:eq:Qsum} therefore
suggests $e^{-2\rho}$, but uniform summation and all-order error
control must still be justified.

\subsection{An explicit algorithm for every coefficient}
Define the following polynomial in $j$, with coefficients in
$\Q[\rho^{-1}]$:
\begin{equation}
 L_\ell(j;\rho)=\frac{(-1)^{\ell+1}}{\ell}
 \left\{\sum_{a=0}^{j-1}\left(\frac{2j-2-a}{2}\right)^\ell
       +\sum_{a=0}^{j-1}\left(-\frac a\rho\right)^\ell
       -2\sum_{a=1}^{j}a^\ell\right\}.
 \label{lis:eq:Lell}
\end{equation}
Faulhaber's formulas interpret each finite power sum polynomially.
Set
\begin{equation}
 P_0(j;\rho)=1,\qquad
 P_d(j;\rho)=\frac1d\sum_{\ell=1}^{d}
                    \ell L_\ell(j;\rho)P_{d-\ell}(j;\rho).
 \label{lis:eq:Pell}
\end{equation}
This is the coefficient recurrence for exponentiating
$\sum_{\ell\ge1}L_\ell u^\ell$.
Finally, define
\begin{equation}
 c_d(\rho)=e^{2\rho}\sum_{j=0}^{\infty}
                  \frac{(-2\rho)^j}{j!}P_d(j;\rho).
 \label{lis:eq:cdsum}
\end{equation}
The last expression is elementary, not an unevaluated infinite-sum
obstruction. If $P_d(j;\rho)=\sum_b p_{d,b}(\rho)j^b$, then
\begin{equation}
 c_d(\rho)=\sum_b p_{d,b}(\rho)\,\mathcal T_b(-2\rho),
 \qquad
 \mathcal T_b(x)=\sum_{a=0}^{b}
       \left\{\begin{matrix}b\\a\end{matrix}\right\}x^a.
 \label{lis:eq:Touchard}
\end{equation}
Here the braces denote Stirling numbers of the second kind.
The identity follows by expanding $j^b$ in falling factorials and
summing $\sum_jx^j\fall{j}{a}/j!=x^ae^x$.
No positive-probability interpretation is imposed on the negative
argument $-2\rho$.

For example,
\begin{equation}
 L_1(j;\rho)=-\frac{j^2+7j}{4}
                     -\frac{j(j-1)}{2\rho},\label{lis:eq:L1}
\end{equation}
and $P_1=L_1$. Substitution in \eqref{lis:eq:Touchard} gives
$c_1=2\rho-\rho^2$. Equations~\eqref{lis:eq:Lell}--\eqref{lis:eq:Touchard}
produce all the subsequent coefficients by exact rational arithmetic.

\begin{proposition}[Polynomial structure]\label{lis:prop:polynomial}
For every $d\ge1$, $c_d\in\Q[\rho]$, $c_d(0)=0$, and
\begin{equation}
 \deg c_d=2d,\qquad
 [\rho^{2d}]c_d=\frac{(-1)^d}{d!}.
 \label{lis:eq:degree}
\end{equation}
The continuous value at $\rho=0$ is a polynomial fact; it does not
by itself establish uniformity of the asymptotic expansion there.
\end{proposition}
\begin{proof}
Since $\deg_j L_\ell\le\ell+1$, \eqref{lis:eq:Pell} gives
$\deg_j P_d\le2d$. The only negative powers of $\rho$ in
\eqref{lis:eq:Qproduct} come from $\prod_{a=0}^{j-1}(1-a u/\rho)$,
where $u=1/k$. The coefficient of $\rho^{-b}$ for $b\ge1$ vanishes
at each integer $j=0,1,\ldots,b$: there are fewer than $b$ nonzero
factors $a$ to select. Thus, as a polynomial in $j$, that coefficient
has a falling-factorial expansion involving only $\fall{j}{a}$ with
$a\ge b+1$.

In \eqref{lis:eq:cdsum}, $\rho^{-b}\fall{j}{a}$ becomes
$\rho^{-b}(-2\rho)^a$, which is polynomial and divisible by $\rho$.
For $b=0$ and $d\ge1$, the coefficient vanishes at $j=0$ because the
empty product is identically one, and the same divisibility follows.
The degree bound is $2d$.

To obtain degree $2d$, the only term of degree $2d$ in $j$ with no
negative power of $\rho$ is $(-j^2/4)^d/d!$ from $L_1^d/d!$.
Every $L_\ell$ with $\ell>1$ has degree strictly less than $2\ell$.
Replacing $j^{2d}$ by its Touchard polynomial contributes the leading
term $(-2\rho)^{2d}$. The claimed leading coefficient follows.
\end{proof}

\subsection{Uniform remainder estimates}
\begin{proof}[Proof of Theorem~\ref{lis:thm:asymptotic}]
Fix $0<\rho_0\le\rho\le\rho_1<\infty$.
First, the absolute value of the $j$th summand of
\eqref{lis:eq:normalizedsum} is bounded by
\begin{equation}
 \binom{2k+2j-2}{j}\frac{\fall{m}{j}}{\bigl(\rise{k+1}{j}\bigr)^2}
 \le\frac{[2\rho_1(1+\rho_1)]^j}{j!},\qquad 0\le j\le m.
 \label{lis:eq:domination}
\end{equation}
Indeed, $j\le m\le\rho_1k$, the top of the binomial is at most
$2(1+\rho_1)k$, $\fall{m}{j}\le(\rho_1k)^j$, and
$\rise{k+1}{j}\ge k^j$. This is a summable bound independent of $k$.

For a quantitative expansion, split the sum at $J_k=\lfloor k^{1/4}\rfloor$.
For $j\le J_k$, regard \eqref{lis:eq:Qproduct} as a function $Q_j(u;\rho)$
of $u$ on $0\le u\le1/k$. All logarithmic factors stay away from zero
uniformly. Its logarithm has Taylor coefficients \eqref{lis:eq:Lell};
its $\ell$th derivative is bounded by $C_{I,\ell}(1+j)^{\ell+1}$.
Also $|\log Q_j(u;\rho)|\le C_Ij^2/k$, so $Q_j$ itself is bounded
uniformly on this interval.

Repeated differentiation of an exponential expresses its derivatives
as finite sums of products of these logarithmic derivatives.
Consequently its $(M+1)$st derivative is bounded by
$C_{I,M}(1+j)^{2M+2}$. Taylor's theorem gives
\begin{equation}
 Q_{k,j}(\rho)=\sum_{d=0}^{M}\frac{P_d(j;\rho)}{k^d}
    +O_{I,M}\left(\frac{(1+j)^{2M+2}}{k^{M+1}}\right),
 \quad j\le J_k.\label{lis:eq:Qremainder}
\end{equation}
After multiplication by $(2\rho)^j/j!$, the remainders sum to
$O_{I,M}(k^{-M-1})$ because factorial denominators dominate every
fixed power of $j$.

For $j>J_k$, \eqref{lis:eq:domination} bounds the original summands.
The polynomial approximants have the same type of bound with an extra
fixed power of $j$. For any fixed $C,D,L$, the tail
$\sum_{j>k^{1/4}}C^j(1+j)^D/j!$ is $o(k^{-L})$.
This follows, for example, from $j!\ge(j/e)^j$ followed by a geometric
tail bound once $j$ is sufficiently large. Thus one can extend each
coefficient sum to infinity with an error smaller than every algebraic
order.

Equations~\eqref{lis:eq:Qsum}, \eqref{lis:eq:Qremainder}, and
\eqref{lis:eq:cdsum} now give \eqref{lis:eq:allorders} with $C_1$ in place of
$T$. Proposition~\ref{lis:prop:errorbound} shows that replacing $C_1$ by
$T$ does not change any coefficient or the stated remainder.
\end{proof}

\section{Exact-length counts, overshoots, and subsequence multiplicity}
The uniform tail expansion also gives a local exact-length expansion.
For fixed $N$ and $k$, replacing $k$ by $k+1$ changes the parameters to
\begin{equation}
 \rho' = \frac{\rho k-1}{k+1},\qquad
 \frac{B(N,k+1)}{B(N,k)}=\frac{\rho k}{(k+1)^2}.
 \label{lis:eq:shiftparameters}
\end{equation}
Both ratios admit expansions to arbitrary order in $1/k$ on the same
compact $\rho$ intervals. Subtraction via \eqref{lis:eq:difference} yields:

\begin{corollary}[Exact-length expansion]\label{lis:cor:point}
Uniformly in the regime of Theorem~\ref{lis:thm:asymptotic},
\begin{align}
 \frac{A(N,k)}{B(N,k)}=e^{-2\rho}\bigg[1
 &+\frac{\rho-\rho^2}{k}
 +\frac{\rho^4/2-7\rho^3/3+3\rho^2+\rho/3}{k^2}\notag\\
 &+\frac{-\rho^6/6+11\rho^5/6-41\rho^4/6
            +31\rho^3/3-2\rho^2/3-2\rho}{k^3}
 +O_I(k^{-4})\bigg].\label{lis:eq:pointasymptotic}
\end{align}
All higher coefficients follow by the same substitution and subtraction.
\end{corollary}
\begin{proof}
Apply \eqref{lis:eq:allorders} to $(N,k)$ and $(N,k+1)$, use
\eqref{lis:eq:shiftparameters}, and expand in $1/k$. The remainders stay
uniform because $\rho'$ remains in a slightly enlarged compact interval
inside $(0,\infty)$. Collecting terms gives the displayed polynomials.
\end{proof}

\begin{corollary}[Conditional overshoot]\label{lis:cor:overshoot}
For every fixed integer $d\ge1$, under a uniformly random permutation
of $N$ letters,
\begin{equation}
 \Prob\bigl(\LIS\ge k+d\mid\LIS\ge k\bigr)
 =\left(\frac{\rho}{k}\right)^d\left(1+O_{I,d}(k^{-1})\right).
 \label{lis:eq:overshoot}
\end{equation}
In particular, the conditional probability that $\LIS=k$ is
$1-\rho/k+O_I(k^{-2})$.
\end{corollary}
\begin{proof}
The conditional probability is $T(N,k+d)/T(N,k)$. The ratio of its
$B$ factors is
$\fall{N-k}{d}/\bigl(\rise{k+1}{d}\bigr)^2$, which equals
$(\rho/k)^d(1+O_{I,d}(k^{-1}))$. The exponential factors and the
relative correction series have ratio $1+O_{I,d}(k^{-1})$.
\end{proof}

\begin{corollary}[Mean number of witnesses]\label{lis:cor:witnesses}
Let $Z_k(\pi)$ be the number of increasing subsequences of length $k$.
Then
\begin{equation}
 \E[Z_k\mid\LIS\ge k]
 =e^{2\rho}\left(1-\frac{2\rho-\rho^2}{k}
                         +O_I(k^{-2})\right).
 \label{lis:eq:witnessmean}
\end{equation}
\end{corollary}
\begin{proof}
Equation~\eqref{lis:eq:Bwitness} gives
$\E[Z_k\mid\LIS\ge k]=B/T$. Invert \eqref{lis:eq:allorders}.
\end{proof}

This conclusion concerns the mean only. It does not imply a Poisson
limit, concentration, or convergence of the entire witness-count
distribution. Those require joint witness information not contained in
the first mean.

\section{OEIS diagonals and numerical checks}
\subsection{Five correction terms for A269021}
Set $N=2n$, $k=n$, so $\rho=1$. The coefficient rule gives
\begin{equation}
 (c_0(1),\ldots,c_5(1))
 =\left(1,1,\frac92,7,-\frac{149}{8},-\frac{3911}{40}\right).
 \label{lis:eq:cone}
\end{equation}
Stirling's expansion gives
\begin{equation}
 \frac{B(2n,n)}{16^n(n-1)!/\pi}
 =1-\frac1{4n}+\frac1{32n^2}+\frac1{128n^3}
       -\frac5{2048n^4}-\frac{23}{8192n^5}+O(n^{-6}).
 \label{lis:eq:central}
\end{equation}
One way to check it is to exponentiate
$-1/(4n)+1/(96n^3)-1/(320n^5)+O(n^{-7})$.

\begin{corollary}[A269021]\label{lis:cor:269}
For $a_n=T(2n,n)$,
\begin{align}
 a_n=\frac{16^n(n-1)!}{\pi e^2}\bigg[1
 &+\frac{3}{4n}+\frac{137}{32n^2}+\frac{757}{128n^3}\notag\\
 &-\frac{41429}{2048n^4}
       -\frac{3803959}{40960n^5}+O(n^{-6})\bigg].
 \label{lis:eq:269asymptotic}
\end{align}
The expansion continues to every fixed order by
\eqref{lis:eq:Lell}--\eqref{lis:eq:Touchard} and Stirling's series.
\end{corollary}
\begin{proof}
Multiply \eqref{lis:eq:cone}, interpreted as its $1/n$ expansion from
Theorem~\ref{lis:thm:asymptotic}, by \eqref{lis:eq:central}.
\end{proof}

For the exact-length diagonal $b_n=A(2n,n)$, the same calculation using
\eqref{lis:eq:pointasymptotic} gives
\begin{equation}
 b_n=\frac{16^n(n-1)!}{\pi e^2}
       \left(1-\frac1{4n}+\frac{49}{32n^2}
                              +\frac{273}{128n^3}+O(n^{-4})\right).
 \label{lis:eq:267asymptotic}
\end{equation}
The different $1/n$ coefficients distinguish ``at least $n$'' from
``exactly $n$'', despite their common leading equivalent.
The latter leading equivalent is already proved in \cite{Moews}.

\subsection{Independent exact comparisons}
The supplied verifier uses the hook-length formula and integer
partitions to compute $A$, $T$, and $C_1,\ldots,C_4$ directly. A separate
calculation uses truncated rational power series in the squared-minor
formula \eqref{lis:eq:squaredminors}. It checks genuinely nonzero $j=2$
and $j=3$ cases, not just empty supports.

\begin{center}
\begin{tabular}{@{}r r r@{}}
\toprule
$n$ & $T(2n,n)$ & $A(2n,n)$\\
\midrule
0 & 1 & 1\\
1 & 2 & 1\\
2 & 23 & 13\\
3 & 588 & 381\\
4 & 24\,553 & 17\,557\\
5 & 1\,438\,112 & 1\,100\,902\\
6 & 108\,469\,917 & 87\,116\,283\\
7 & 9\,996\,042\,284 & 8\,312\,317\,976\\
8 & 1\,086\,997\,811\,325 & 927\,716\,186\,325\\
\bottomrule
\end{tabular}
\end{center}
The initial values agree with the corresponding OEIS entries
\cite{oeis269,oeis267}; the small-index empty-permutation convention
is imposed separately from the positive-$k$ array.

The completed exact test run covers $0\le N\le26$, $1\le k\le N+1$:
378 cells. It verifies Plancherel normalization, the first-moment
single sum, the residual identity, the sharp cutoff, and the first
rectangle discrepancy for truncation ranks $1$ through $4$, whenever
the relevant point lies in that range. There are 140 independent
Bessel-minor comparisons. The A269021 formula matches 16 displayed
OEIS terms and generates terms through $n=60$. Symbolic computation
produces the coefficient polynomials through order five and checks
15 finite-product specializations by a separate rational-series method.

\subsection{Accuracy of the asymptotic expansion}
Let $L_n=16^n(n-1)!/(\pi e^2)$, and let $S_M(n)$ retain terms through
$n^{-M}$ in the square brackets of \eqref{lis:eq:269asymptotic}.
The next table reports $|L_nS_M(n)/a_n-1|$.
The integer $a_n$ is evaluated exactly; the rescaling uses 100 decimal
digits of floating-point precision.
\begin{center}
\begin{tabular}{@{}r r r r r@{}}
\toprule
$n$ & $M=0$ & $M=1$ & $M=3$ & $M=5$\\
\midrule
20 & $4.65\cdot10^{-2}$ & $1.08\cdot10^{-2}$ & $1.44\cdot10^{-4}$ & $4.38\cdot10^{-6}$\\
50 & $1.65\cdot10^{-2}$ & $1.73\cdot10^{-3}$ & $3.46\cdot10^{-6}$ & $1.62\cdot10^{-8}$\\
100 & $7.87\cdot10^{-3}$ & $4.30\cdot10^{-4}$ & $2.10\cdot10^{-7}$ & $2.40\cdot10^{-10}$\\
200 & $3.84\cdot10^{-3}$ & $1.07\cdot10^{-4}$ & $1.29\cdot10^{-8}$ & $3.64\cdot10^{-12}$\\
400 & $1.90\cdot10^{-3}$ & $2.68\cdot10^{-5}$ & $7.98\cdot10^{-10}$ & $5.60\cdot10^{-14}$\\
\bottomrule
\end{tabular}
\end{center}
These floating-point comparisons are diagnostic, not interval-certified
error bounds. The proof of the asymptotic statement is the uniform
remainder argument, not agreement in the table. The result is an
asymptotic expansion: adding terms need not improve an approximation
at every small $n$ or at arbitrarily high truncation order.

\section{Further research questions and topics}
\label{lis:sec:further}
The following questions are separate from the theorems proved above.
They are ordered roughly by proximity to the present construction,
not by a claim of difficulty or priority.

\begin{enumerate}[leftmargin=*,itemsep=8pt]
\item \textbf{Certify the specific one-third-sector recurrences.}
Use the $C_1-C_2$ representation to construct certificates for the
operators guessed in \cite{KW}. The existence theorem does not itself
identify those operators. A complete result would publish the
certificates, initial region, and singular-index treatment.

\item \textbf{Determine recurrence complexity as the sector widens.}
Find upper and lower bounds on minimal pure-shift orders, degrees, and
holonomic ranks as functions of $R$. The fixed-dimensional multisum
proves finiteness, but the naive determinant expansion has $(R!)^2$
outer terms at its largest moment. Can a structured representation
avoid that growth?

\item \textbf{Compute higher moments without full determinant expansion.}
Exploit Cauchy--Binet, symmetric-polynomial coordinates, or discrete
orthogonality to obtain efficient exact formulas for $C_2,C_3$, and
beyond. Any claimed speedup should distinguish arithmetic-operation
counts from bit complexity and should be compared with independent
tableau counts.

\item \textbf{Resolve the exponentially small sectors.}
Find the first asymptotic equivalent of $C_1-T$ in a region with
$N>2k+2$, and then of the residual after $R$ moments. The bound
\eqref{lis:eq:superalgebraic} does not determine the exponential action,
prefactor, or all-order corrections. Near a rectangle boundary, a
fixed excess number of boxes may offer a tractable starting regime.

\item \textbf{Extend uniformity beyond compact linear regimes.}
Investigate $\rho\to0$ with growing deficit, and $\rho\to\infty$ with
$k=o(N)$. These are not covered by Theorem~\ref{lis:thm:asymptotic}.
The polynomial coefficients alone do not justify substitution of a
growing $\rho$; one needs new remainder estimates and, eventually,
a different balance of moment contributions.

\item \textbf{Study the conditional witness distribution.}
The mean number of length-$k$ increasing subsequences tends to
$e^{2\rho}$. Determine higher conditional moments, limiting laws,
and structural descriptions of permutations in the rare tail event.
Shifted-row factorial moments and witness-count factorial moments
must be kept distinct.

\item \textbf{Classify rational-ray OEIS families.}
Generate exact values and asymptotic expansions for
$T(an+a_0,bn+b_0)$ and $A(an+a_0,bn+b_0)$, and identify which already
have OEIS entries. The rational-ray theorem supplies a recurrence
existence guarantee before guessing begins. Each proposed entry or
formula should receive its own source and duplicate audit.

\item \textbf{Formalize the finite proof boundary.}
A staged formalization could first verify the rectangle lemma,
binomial residual, exact finite coefficient identity, and low-rank
certificates. A full formalization must additionally connect
Robinson--Schensted, the Bessel correlation theorem, and holonomic
closure to the precise array semantics. Importing any of these as
an axiom should be explicitly recorded rather than hidden by a
finite numerical checker.
\end{enumerate}

\section{Conclusion and proof audit}
The fixed-sector theorem follows from a short conceptual chain:
\[
 \text{shifted-row threshold}
 \ \Longrightarrow\ \text{rectangle obstruction}
 \ \Longrightarrow\ \text{finite inclusion--exclusion},
\]
followed by the Bessel-kernel conversion of every fixed moment into a
proper-hypergeometric multisum. The number of summation variables is
fixed before $N,k$ vary. This is the precise point that allows
holonomic closure and establishes all fixed sectors.

The first moment has a much simpler single-sum expression. Its
normalized summands converge to a signed exponential series, while
factorial domination makes the expansion uniform to every algebraic
order. The higher shifted-row contributions are either forbidden
exactly or exponentially smaller than those orders. This gives a
single coefficient calculus for the full positive-linear regime,
including the two OEIS diagonals and their local probability
consequences.

For audit purposes, the proof has four distinct layers. The external
inputs are established named theorems, cited where used. The algebraic
reductions are displayed explicitly, including the factorial arguments
and support conditions needed by holonomic summation. The asymptotic
claims include uniform remainder bounds. The executable tests provide
independent finite checks, not a replacement for any infinite theorem.
No specific high-order guessed recurrence, global holonomicity of the
original array, or formal proof-assistant certification is claimed.

\appendix
\section{Reproduction guide and package contents}
The package is self-contained for its computations and does not need
network access. It contains the article source and compiled PDF,
\texttt{verify.py}, a requirements file, a source-audit note, and
machine-readable results. Run
\begin{lstlisting}
python verify.py --output results
\end{lstlisting}
The exact combinatorial checks need only Python's standard library.
The symbolic and high-precision parts use SymPy and mpmath. To run
only the exact checks, use
\begin{lstlisting}
python verify.py --core-only --output results
\end{lstlisting}
The source can be compiled with
\begin{lstlisting}
pdflatex -interaction=nonstopmode -halt-on-error article.tex
pdflatex -interaction=nonstopmode -halt-on-error article.tex
\end{lstlisting}

The principal output files are
\begin{center}
\begin{tabular}{@{}p{0.38\textwidth}p{0.56\textwidth}@{}}
\toprule
File under \texttt{results/} & Contents\\
\midrule
\texttt{small\_counts.csv} & Exact-length counts, tails, and four moments for small $N,k$.\\
\texttt{bessel\_checks.csv} & Independent rational-series evaluations of fixed moments.\\
\texttt{a269021.csv}, \texttt{a267433.csv} & Exact tail and exact-length diagonal values for $0\le n\le60$.\\
\texttt{symbolic\_coefficients.json} & The five computed polynomial corrections and diagonal specializations.\\
\texttt{asymptotic\_checks.csv} & Signed relative errors of four asymptotic truncations.\\
\texttt{uniform\_checks.csv} & Off-diagonal tests at $\rho=1/2,1,2$.\\
\texttt{verification\_summary.json} & Test scope and pass status.\\
\bottomrule
\end{tabular}
\end{center}

The Bessel verifier factors
$J_d(2\sqrt t)=t^{d/2}F_d(t)$, where
$F_d(t)=\sum_s(-1)^st^s/(s!(d+s)!)$. For each increasing pair $u,v$
with budget \eqref{lis:eq:basedegree}, it forms $\det F$, squares it,
and convolves with $e^t$. All of this is done with exact rational
coefficients. The independent tableau computation uses no Bessel
functions. This separation detects normalization, shift, missing
factorial, and determinant-sign errors that would be difficult to
spot from a single formula alone.

\begin{writenote}[Added 3 October 2026, batch~85]
This appendix describes the delivered package. In ProveIt the verifier
is \texttt{code/verify.py}, the files listed under \texttt{results/} are
shipped in \texttt{data/}, beside \texttt{data/requirements.txt}, and
\texttt{build.sh} is \texttt{code/build.sh} (it runs three
\texttt{pdflatex} passes, not the two shown above); the source-audit note
is \texttt{SOURCE\_AUDIT.md}. The delivered PDF is not shipped; the
committed \texttt{article.pdf} is a rebuild including these notes. Run
from \texttt{code/} without \texttt{-{}-output}, the verifier would write
a new \texttt{code/results/} directory; the report's README gives a rerun
on a scratch copy that never writes into the report. The sixteen terms
each of A269021 and A267433 embedded in the verifier are OEIS data, under
CC~BY-SA~4.0. At intake (3 October 2026) the full run passed on such a
copy in 71~s (the delivered summary records 4.6~s on the producer's
machine); every regenerated data file matched the shipped one, the two
text-mode outputs up to line endings, and the summary differed only in
its elapsed time.
\end{writenote}

\section{Additional coefficient polynomials}
The fourth and fifth uniform tail coefficients, generated by the
all-orders rule and checked by the supplied symbolic code, are
\begin{align}
 c_4(\rho)=\frac{\rho}{360}\big(&15\rho^7-360\rho^6+3440\rho^5
 -16716\rho^4+38960\rho^3\notag\\
 &-19080\rho^2-15100\rho+2136\big),\label{lis:eq:c4}\\
 c_5(\rho)=-\frac{\rho}{360}\big(&3\rho^9-110\rho^8+1670\rho^7
 -13796\rho^6+66352\rho^5\notag\\
 &-170464\rho^4+145380\rho^3+67016\rho^2
 -54612\rho-6240\big).\label{lis:eq:c5}
\end{align}
These expressions concern the normalization $B(N,k)e^{-2\rho}$.
They should not be substituted directly into the A269021 normalization
$16^n(n-1)!/(\pi e^2)$ without also applying the central-binomial
correction \eqref{lis:eq:central}.

\begin{thebibliography}{99}
\addcontentsline{toc}{section}{References}
\bibitem{KW}
M. Kauers and C. Wang,
\emph{Recurrences for permutations with long increasing subsequences},
arXiv:2609.02220v1, September 2, 2026.
\url{https://arxiv.org/abs/2609.02220}.

\bibitem{KK}
M. Kauers and C. Koutschan,
\emph{Some D-finite and some possibly D-finite sequences in the OEIS},
Journal of Integer Sequences \textbf{26} (2023), Article 23.4.5.
\url{https://arxiv.org/abs/2303.02793}.

\bibitem{BOO}
A. Borodin, A. Okounkov, and G. Olshanski,
\emph{Asymptotics of Plancherel measures for symmetric groups},
Journal of the American Mathematical Society \textbf{13} (2000),
481--515. Theorem 2 and Proposition 2.9 in the preprint version.
\url{https://arxiv.org/abs/math/9905032}.

\bibitem{WZ}
H. S. Wilf and D. Zeilberger,
\emph{Rational function certification of multisum/integral/``$q$'' identities},
Bulletin of the American Mathematical Society (N.S.) \textbf{27} (1992),
148--153.
\url{https://arxiv.org/abs/math/9207218}.

\bibitem{Koutschan}
C. Koutschan,
\emph{Creative telescoping for holonomic functions},
in C. Schneider and J. Bl\"umlein (eds.),
\emph{Computer Algebra in Quantum Field Theory}, Springer, 2013,
171--194.
\url{https://arxiv.org/abs/1307.4554}.

\bibitem{Schensted}
C. Schensted,
\emph{Longest increasing and decreasing subsequences},
Canadian Journal of Mathematics \textbf{13} (1961), 179--191.
\url{https://doi.org/10.4153/CJM-1961-015-3}.

\bibitem{Moews}
D. Moews,
Answer to \emph{Reference or proof for number of permutations of [2n]
with longest increasing subsequence of length n},
Mathematics Stack Exchange, September 22, 2023, revised September 25, 2023.
\url{https://math.stackexchange.com/questions/4771623}.

\bibitem{oeisA}
OEIS Foundation, \emph{A047874: permutations by exact longest increasing
subsequence length}. \url{https://oeis.org/A047874}.

\bibitem{oeisT}
OEIS Foundation, \emph{A214152: permutations containing an increasing
subsequence of a specified length}. \url{https://oeis.org/A214152}.

\bibitem{oeis269}
OEIS Foundation, \emph{A269021: permutations of [2n] containing an
increasing subsequence of length n}. \url{https://oeis.org/A269021}.

\bibitem{oeis267}
OEIS Foundation, \emph{A267433: permutations of [2n] with longest
increasing subsequence of length n}. \url{https://oeis.org/A267433}.

\bibitem{repo}
V. Reshetnikov, \emph{ProveIt}, GitHub repository,
revision {\small\texttt{6bf7f30d0352f7596e70928b3d4f304914075907}}.
\url{https://github.com/VladimirReshetnikov/ProveIt}.
\end{thebibliography}
\noindent\small Web and repository source audit: October 3, 2026.
OEIS entries are living documents; the audit records the versions inspected,
not an assertion about all subsequent updates.
\end{document}
